\documentclass[11pt]{article}
\usepackage{graphicx} % Required for inserting images
\usepackage{booktabs} % \toprule / \midrule / \bottomrule in tables
\usepackage{amsmath}
\usepackage{tabularx}
\usepackage[section]{placeins}
\usepackage[margin=1in]{geometry}
\usepackage{microtype}
\usepackage{float}
\usepackage[font=small,labelfont=bf]{caption}

\raggedbottom
\renewcommand{\arraystretch}{1.15}
\setlength{\tabcolsep}{5pt}
\setlength{\intextsep}{12pt plus 2pt minus 2pt}
\captionsetup[table]{position=top,skip=6pt}
\captionsetup[figure]{skip=6pt}

\title{Neural Networks --- Final Project}
\author{Noam Major - 208941989 \and Ziv Assoulin - 325312999}
\date{September 2026}

\begin{document}

\maketitle

\section{Part A}

\noindent\textbf{Spike-Train Analysis (neurodata.npy)}

Data: injected stimulus current (pA) and the voltage responses (mV) of two neurons (spiny, aspiny), sampled at 50 kHz. The aspiny neuron is identified as the one with the higher firing rate.

\FloatBarrier
\subsection{Question 1: Spiny and Aspiny Neurons}

Neurons are nerve cells that receive and transmit information. Their branching extensions, called \textbf{dendrites}, receive signals from other neurons at connections called \textbf{synapses}. A neuron can respond by producing a brief electrical pulse, called an \textbf{action potential} or \textbf{spike}. The terms \emph{spiny} and \emph{aspiny} describe the appearance of its dendrites:

\begin{itemize}
    \item \textbf{Spiny neurons} have many small protrusions called \emph{dendritic spines}, which provide sites for incoming synaptic connections. In the cerebral cortex, these neurons are typically \emph{excitatory}: their signals tend to make other neurons more likely to fire. A common example is the pyramidal cell, which releases the chemical messenger \emph{glutamate}.
    \item \textbf{Aspiny neurons} lack these spines, so their dendrites look smoother, but they still receive synaptic inputs. In the cortex, they are commonly \emph{inhibitory interneurons}: cells that regulate nearby neurons by making them less likely to fire. Examples include basket cells, which release the chemical messenger \emph{GABA}. Many can fire rapidly and help control the timing of activity in a neural circuit.
\end{itemize}

These are typical associations, not universal rules: dendritic appearance alone does not determine a neuron's function.

\FloatBarrier
\subsection{Question 2: Electrical Stimulus over Time}

The stimulus is an injected current with two pulses: a short 50 pA calibration pulse ($\sim$10 ms, 0.005-0.015 s) used to estimate input resistance, followed by the main 150 pA, 1-second current step (1.02-2.02 s) that drives repetitive spiking. Current returns to 0 pA afterward.

\begin{figure}[H]
\centering
\includegraphics[width=0.85\linewidth,height=0.23\textheight,keepaspectratio]{figures/parts-a-b/figure-a2.png}
\caption{Injected stimulus current (pA) vs. time (s).}
\label{fig:a2}
\end{figure}

\FloatBarrier
\subsection{Question 3: Electrical Activity of Spiny and Aspiny Neurons}

Both neurons fire action potentials during the current step. The main difference is firing rate: the aspiny neuron fires much more frequently (91 spikes) than the spiny neuron (25 spikes) over the same 1-second step, consistent with the fast-spiking phenotype typical of many inhibitory interneurons.

\begin{figure}[H]
\centering
\includegraphics[width=0.65\linewidth,height=0.27\textheight,keepaspectratio]{figures/parts-a-b/figure-a3.png}
\caption{Membrane voltage (mV) vs. time (s) for the spiny and aspiny neurons.}
\label{fig:a3}
\end{figure}

\FloatBarrier
\subsection{Question 4: Spike Trains and Action-Potential Threshold}

Both neurons rest around -70 to -80 mV and their spikes peak well above 0 mV, so a threshold of -20 mV sits in an empty gap between the two. A sweep over thresholds from -50 to +4 mV confirms the detected spike count is exactly constant (25 for spiny, 91 for aspiny) across roughly -35 to +2 mV, so the exact value inside that range does not matter.

\begin{itemize}
    \item Too low (e.g. -60 mV): falls inside the resting-potential noise band $\rightarrow$ false positive spikes, inflated firing rate, corrupted ISI/CV statistics.
    \item Too high (e.g. +40 mV): exceeds some real spike peaks (aspiny peaks reach only $\sim$22 mV) $\rightarrow$ missed spikes, underestimated firing rate.
\end{itemize}

\begin{figure}[H]
\centering
\includegraphics[width=0.92\linewidth,height=0.22\textheight,keepaspectratio]{figures/parts-a-b/figure-a4a.png}
\caption{Detected spike count vs. threshold sweep, confirming a stable plateau around the chosen -20 mV threshold.}
\label{fig:a4a}
\end{figure}

\begin{figure}[H]
\centering
\includegraphics[width=0.85\linewidth,height=0.22\textheight,keepaspectratio]{figures/parts-a-b/figure-a4b.png}
\caption{Spike-train raster plots for the spiny and aspiny neurons.}
\label{fig:a4b}
\end{figure}

\begin{figure}[H]
\centering
\includegraphics[width=0.78\linewidth,height=0.28\textheight,keepaspectratio]{figures/parts-a-b/figure-a4c.png}
\caption{Zoomed voltage trace with the -20 mV detection threshold and detected spikes marked.}
\label{fig:a4c}
\end{figure}

\FloatBarrier
\subsection{Question 5: Inter-Spike Intervals (ISI) and Coefficient of Variation (CV)}

\begin{table}[H]
\centering
\caption{Inter-spike intervals and firing regularity during the current step.}
\label{tab:a5}
\small
\begin{tabularx}{\textwidth}{@{}lccc>{\raggedright\arraybackslash}X@{}}
\toprule
\textbf{Neuron} & \textbf{CV} & \textbf{ISI range (ms)} & \textbf{n intervals} & \textbf{Firing pattern} \\
\midrule
Spiny & 0.070 & 29.9 -- 44.8 & 24 & Pacemaker-like (CV $\ll$ 1) \\
Aspiny & 0.049 & 9.7 -- 12.8 & 90 & Pacemaker-like (CV $\ll$ 1) \\
\bottomrule
\end{tabularx}
\end{table}

Both CVs are far below 1, so both neurons fire highly regularly during the stimulus (pacemaker-like, not Poisson). The aspiny neuron is even more regular despite firing $\sim$4x faster. ISIs lengthen slightly and monotonically over the step -- spike-frequency adaptation -- which is the main source of the small variability captured by the CV.

\begin{figure}[H]
\centering
\includegraphics[width=0.92\linewidth,height=0.24\textheight,keepaspectratio]{figures/parts-a-b/figure-a5.png}
\caption{ISI distribution histograms (ms) for the spiny and aspiny neurons.}
\label{fig:a5}
\end{figure}

\FloatBarrier
\subsection{Question 6: Time-Dependent Firing Rate}

Two estimators are shown: a 50 ms-binned PSTH (spike count / bin width) and a Gaussian-kernel estimate (sigma = 50 ms). The kernel is continuous and insensitive to bin placement; sigma = 50 ms is short enough to resolve the onset/offset of the 1 s current step while long enough (relative to the 10-45 ms ISIs) to average over several spikes for a smooth estimate. Both methods agree closely, as expected given how regular (low-CV) the firing is.

\begin{figure}[H]
\centering
\includegraphics[width=0.80\linewidth,height=0.34\textheight,keepaspectratio]{figures/parts-a-b/figure-a6.png}
\caption{Time-dependent firing rate (Hz) vs. time (s): binned PSTH vs. Gaussian-kernel estimate.}
\label{fig:a6}
\end{figure}

\FloatBarrier
\subsection{Question 7: Average Firing Rates and Probability Distributions}

Average firing rate during the stimulus: spiny = 25.00 Hz, aspiny = 91.00 Hz. Gaussian distributions were built around these means with sigma\_spiny = 15 Hz and sigma\_aspiny = 20 Hz, chosen to give the two curves clear but partial overlap (about a third of the 66 Hz gap between the means).

\begin{figure}[H]
\centering
\includegraphics[width=0.70\linewidth,height=0.29\textheight,keepaspectratio]{figures/parts-a-b/figure-a7.png}
\caption{Gaussian probability distributions of firing rate (Hz) for spiny vs. aspiny neurons.}
\label{fig:a7}
\end{figure}

\FloatBarrier
\subsection{Question 8: Neuron Classifier Based on Distributions}

The classification threshold z = 54.59 Hz is chosen at the point where the two Gaussian pdfs cross (equal class-conditional density under equal priors), rather than the plain midpoint of the means. The classifier was evaluated on 1000 synthetic samples drawn from each Gaussian (2000 total).

\begin{table}[H]
\centering
\caption{Performance of the firing-rate classifier on 2,000 synthetic samples.}
\label{tab:a8}
\small
\begin{tabular}{@{}lr@{}}
\toprule
\textbf{Metric} & \textbf{Value} \\
\midrule
Accuracy & 0.974 \\
Precision (Aspiny) & 0.976 \\
Recall (Aspiny) & 0.972 \\
Precision (Spiny) & 0.972 \\
Recall (Spiny) & 0.976 \\
F1 (Aspiny) & 0.974 \\
\bottomrule
\end{tabular}
\end{table}

\begin{figure}[H]
\centering
\begin{minipage}[t]{0.55\linewidth}
\centering
\vspace{0pt}
\includegraphics[width=\linewidth]{figures/parts-a-b/figure-a8a.png}
\caption{Firing-rate distributions with the classification threshold z = 54.59 Hz marked.}
\label{fig:a8a}
\end{minipage}\hfill
\begin{minipage}[t]{0.41\linewidth}
\centering
\vspace{0pt}
\includegraphics[width=\linewidth]{figures/parts-a-b/figure-a8b.png}
\caption{Confusion matrix for the z-threshold classifier (rows: true label, columns: predicted label).}
\label{fig:a8b}
\end{minipage}
\end{figure}

\FloatBarrier
\subsection{Question 9: Likelihood Ratios}

\begin{table}[H]
\centering
\caption{Likelihood ratios and classifications for the five tested firing rates.}
\label{tab:a9}
\small
\begin{tabular}{@{}rrrl@{}}
\toprule
\textbf{Rate (Hz)} & \textbf{log-LR (Aspiny/Spiny)} & \textbf{LR (Aspiny/Spiny)} & \textbf{Classification} \\
\midrule
25.00 & -5.733 & 0.0032 & Spiny \\
41.50 & -2.745 & 0.064 & Spiny \\
58.00 & 0.771 & 2.16 & Aspiny \\
74.50 & 4.817 & 123.6 & Aspiny \\
91.00 & 9.392 & 1.2e4 & Aspiny \\
\bottomrule
\end{tabular}
\end{table}

LR = P(rate \textbar{} Aspiny) / P(rate \textbar{} Spiny); LR $>$ 1 favors Aspiny, LR $<$ 1 favors Spiny. The two rates nearest mu\_spiny classify as Spiny, the two nearest mu\_aspiny classify as Aspiny, and the middle rate (near the Q8 threshold z) sits at the decision boundary.

What could change these values:

\begin{itemize}
    \item A shift in either neuron's true mean firing rate would move the distributions and could flip the LR sign for a fixed rate.
    \item A wider variance (more noise, shorter measurement window) increases overlap and pulls LR toward 1 over a broader range, making the two types harder to distinguish.
    \item Unequal class priors: LR alone ignores how common each neuron type is; the decision should then use the posterior odds (LR x prior ratio), not LR alone.
\end{itemize}

\begin{figure}[H]
\centering
\includegraphics[width=0.70\linewidth,height=0.29\textheight,keepaspectratio]{figures/parts-a-b/figure-a9.png}
\caption{The five tested firing rates marked on the Q7 distributions, colored by classification.}
\label{fig:a9}
\end{figure}

\clearpage
\section{Part B}

\noindent\textbf{Unsupervised Learning (feat\_data.csv)}

Data: 1424 neurons x 41 standardized numeric electrophysiological features. 700 spiny, 724 aspiny.

\FloatBarrier
\subsection{Question 1: PCA}

\subsubsection*{(a) PCA to 3 Dimensions}

\begin{figure}[H]
\centering
\includegraphics[width=0.48\linewidth,height=0.25\textheight,keepaspectratio]{figures/parts-a-b/figure-b1a-1.png}
\par\medskip
\includegraphics[width=0.98\linewidth,height=0.20\textheight,keepaspectratio]{figures/parts-a-b/figure-b1a-2.png}
\caption{PCA projection (PC1, PC2, PC3) colored by neuron type; 3D scatter plus pairwise 2D projections with \% variance explained on each axis.}
\label{fig:b1a}
\end{figure}

\FloatBarrier
\subsubsection*{(b) What Do the PCA Axes Represent?}

Each principal component is a linear combination of the 41 standardized features, ordered by the amount of variance it captures: PC1 captures the most variance, PC2 the most remaining variance orthogonal to PC1, and PC3 the next. Capturing the most variance does not mean an axis is the best indicator for separating spiny from aspiny -- PCA is unsupervised and has no notion of the class label.

\FloatBarrier
\subsubsection*{(c) Do Some Axes Separate the Neuron Types Better?}

\begin{table}[H]
\centering
\caption{Single-component classification performance for the first three principal components.}
\label{tab:b1c}
\small
\begin{tabularx}{\textwidth}{@{}lc>{\centering\arraybackslash}X@{}}
\toprule
\textbf{Axis} & \textbf{\% Variance Explained} & \textbf{5-fold CV accuracy (single-axis classifier)} \\
\midrule
PC1 & 32.7\% & 0.798 \\
PC2 & 16.4\% & 0.640 \\
PC3 & 13.6\% & 0.650 \\
\bottomrule
\end{tabularx}
\end{table}

PC1 separates the two neuron types clearly better than PC2 or PC3: its 1D distribution shows the least overlap, and a logistic-regression classifier trained on PC1 alone reaches $\sim$0.80 cross-validated accuracy vs. $\sim$0.64-0.65 for PC2/PC3. This is because PC1 captures by far the most variance, and the biological spiny/aspiny distinction (spike width, adaptation, input resistance) happens to be one of the dominant sources of variance in this feature set.

\begin{figure}[H]
\centering
\includegraphics[width=0.98\linewidth,height=0.22\textheight,keepaspectratio]{figures/parts-a-b/figure-b1c.png}
\caption{Per-axis 1D density of PC1/PC2/PC3 values, split by neuron type.}
\label{fig:b1c}
\end{figure}

\FloatBarrier
\subsubsection*{(d) PCA Axes Needed for $\leq$1\% Reconstruction MSE}

27 of the 41 components are needed to reach $\geq$99\% cumulative explained variance, equivalently $\leq$1\% mean-squared reconstruction error. Because the data is standardized (each feature has unit variance), PCA reconstruction MSE equals exactly the unexplained variance fraction -- verified directly via \texttt{PCA.\allowbreak inverse\_\allowbreak transform}: k=24 $\rightarrow$ 1.68\% MSE, k=27 $\rightarrow$ 0.97\% MSE, k=30 $\rightarrow$ 0.51\% MSE, each matching the value predicted from 1 - cumulative variance.

\begin{figure}[H]
\centering
\includegraphics[width=0.75\linewidth,height=0.30\textheight,keepaspectratio]{figures/parts-a-b/figure-b1d.png}
\caption{Cumulative explained variance vs. number of PCA components, with the 99\% threshold and the 27-component answer marked.}
\label{fig:b1d}
\end{figure}

\FloatBarrier
\subsection{Question 2: Autoencoder}

Architecture: input(41) $\rightarrow$ hidden(32) $\rightarrow$ bottleneck(3) $\rightarrow$ hidden(32) $\rightarrow$ output(41), 3 hidden layers (32, 3, 32), activation = tanh (from the ID-parity rule: the last digits of our IDs are 9 and 9, which sum to 18, an even number). Three variants were trained and compared: a standard AE, a denoising AE (Gaussian noise, factor 0.3, added to the input only; loss always computed against the clean target), and a weight-decay AE (L2 = 1e-3). Optimizer: Adam, lr = 0.005, batch size 32, 150 epochs. The data was split 80/20 into training and validation sets, and the features were standardized with a scaler fitted on the training rows only, so no validation statistics leak into training.

\FloatBarrier
\subsubsection*{(a) Training and Validation MSE}

\begin{table}[H]
\centering
\caption{Final training and validation reconstruction error for each autoencoder variant.}
\label{tab:b2a}
\small
\begin{tabular}{@{}lrr@{}}
\toprule
\textbf{Variant} & \textbf{Final train MSE} & \textbf{Final val MSE} \\
\midrule
Standard AE & 0.3029 & 0.3627 \\
Denoising AE (noise=0.3) & 0.3133 & 0.3628 \\
Weight-Decay AE (L2=1e-3) & 0.3350 & 0.3758 \\
\bottomrule
\end{tabular}
\end{table}

\begin{figure}[H]
\centering
\includegraphics[width=1.00\linewidth,height=0.24\textheight,keepaspectratio]{figures/parts-a-b/figure-b2a.png}
\caption{Train and validation MSE loss vs. epoch, for each of the three autoencoder variants.}
\label{fig:b2a}
\end{figure}

Train and validation MSE stay close together for all three variants (final gap $\sim$0.04-0.06), so none shows strong overfitting given the model's small capacity (2,924 parameters) vs. the 1,139-sample training set. The Denoising and Weight-Decay AEs converge to a slightly higher train MSE than the Standard AE, as expected -- both are deliberately regularized against fitting the training data too closely. Whether that regularization also gives a more class-relevant bottleneck is tested in Q2.c.

\FloatBarrier
\subsubsection*{(b) Low-Dimensional (Bottleneck) Representation}

\begin{figure}[H]
\centering
\includegraphics[width=1.00\linewidth,height=0.27\textheight,keepaspectratio]{figures/parts-a-b/figure-b2b.png}
\caption{3D bottleneck scatter plots (Latent 1-3) for each autoencoder variant, colored by neuron type.}
\label{fig:b2b}
\end{figure}

\FloatBarrier
\subsubsection*{(c) Autoencoder vs. PCA (Qualitative \& Quantitative)}

\begin{table}[H]
\centering
\caption{Class separation and cross-validated classifier accuracy in the three-dimensional representations.}
\label{tab:b2c}
\small
\begin{tabular}{@{}lcc@{}}
\toprule
\textbf{Representation} & \textbf{Silhouette score} & \textbf{5-fold CV accuracy (LogReg)} \\
\midrule
PCA-3 & 0.2276 & 0.8680 \\
Standard AE & 0.3037 & 0.8947 \\
Denoising AE & 0.2989 & 0.8975 \\
Weight-Decay AE & 0.2795 & 0.8778 \\
\bottomrule
\end{tabular}
\end{table}

Qualitatively, the AE bottlenecks show the two classes forming slightly more distinct blobs than the rigid, linear PCA projection. Quantitatively, all three AE variants beat PCA-3 on both silhouette score and cross-validated linear-classifier accuracy. Among the AE variants the two metrics do not agree on a single winner, and the top two are nearly tied: the Standard AE has the highest silhouette (0.304 vs.\ 0.299), while the Denoising AE has the highest cross-validated accuracy (0.898 vs.\ 0.895). We select by cross-validated accuracy, so the Denoising AE is carried forward to Q2.d/Q2.e.

\FloatBarrier
\subsubsection*{(d) Reconstruction Examples}

\begin{figure}[H]
\centering
\includegraphics[width=1.00\linewidth,height=0.34\textheight,keepaspectratio]{figures/parts-a-b/figure-b2d.png}
\caption{Original vs. reconstructed feature values (z-scores) for the 3 best and 3 worst validation-set reconstructions (Denoising AE).}
\label{fig:b2d}
\end{figure}

\begin{sloppypar}
The worst reconstructions are dominated by a single feature, \texttt{latency}, which accounts for almost all of their squared error (mean squared error $\sim$111, against $\sim$2-4 for the next features: \texttt{threshold\_v\_long\_square}, \texttt{threshold\_v\_ramp}, \texttt{trough\_i}, \texttt{fast\_trough\_i} and \texttt{upstroke\_downstroke\_ratio}). The worst-reconstructed samples also tend to have a higher average \textbar{}z-score\textbar{} across all 41 features (0.94, 2.32, 1.44) than the best-reconstructed samples (0.78, 0.48, 1.01) or the validation-set average (0.81) -- i.e., they are mostly statistical outliers. With only 3 latent dimensions to summarize 41 features, the network represents common, population-typical feature combinations well but regresses atypical neurons toward the "average" shape it has learned, losing exactly what makes them unusual.
\end{sloppypar}

\FloatBarrier
\subsubsection*{(e) Did the Autoencoder Learn a Better Representation?}

Yes, modestly but consistently. All three AE variants beat PCA-3 on both silhouette score (0.228 $\rightarrow$ 0.28-0.30) and 5-fold cross-validated classifier accuracy (0.868 $\rightarrow$ 0.88-0.90); the selected variant, the Denoising AE, raises accuracy to 0.898 and silhouette to 0.299. The gain is real but not dramatic: PCA already does well (0.868) because a single linear axis (PC1) already captures much of the spiny/aspiny distinction (Q1.c), leaving limited room for a nonlinear method to improve further. Regularization made little difference to class separation: the Denoising AE and the Standard AE are nearly tied (accuracy 0.898 vs.\ 0.895, silhouette 0.299 vs.\ 0.304), and the Weight-Decay AE is the lowest of the three on both metrics (0.878, 0.280).

\clearpage

\section{Part C}

\paragraph{Setup (common to all questions).}
The feature table (1{,}424 cells, 41 numeric electrophysiology features, label \texttt{dendrite\_type}) was split 80/20 with a stratified random split (\texttt{random\_state=42}): 1{,}139 training cells and 285 test cells (145 aspiny / 140 spiny). The label column was never used as an input feature. All features were standardized with a \texttt{StandardScaler} fitted on the training set only and applied to the test set, so no test statistics leak into preprocessing. Throughout, \emph{spiny} is the positive class.

\FloatBarrier
\subsection{Question 1}

\FloatBarrier
\subsubsection*{(a) Error analysis of the logistic regression}
The main results are: accuracy $=0.9860$, precision $=0.9928$, recall $=0.9786$, F1 $=0.9856$.
The confusion matrix (Figure~\ref{fig:conf_log}) gives: true negatives 144, false positives 1, false negatives 3, true positives 137.

The model makes 4 errors out of 285 samples, but they are not symmetric:
\begin{itemize}
    \item 3 false negatives --- spiny neurons classified as aspiny
    \item 1 false positive --- an aspiny neuron classified as spiny
\end{itemize}

This is why recall (0.9786) sits below precision (0.9928): the model is somewhat
conservative about assigning the spiny label. Almost everything it calls spiny really
is spiny, but it misses three genuine spiny neurons. Biologically, this suggests that a small
subset of spiny neurons carries features close enough to the aspiny profile
that they fall on the wrong side of the decision boundary.

\begin{figure}[H]
\centering
\includegraphics[width=0.45\linewidth,height=0.36\textheight,keepaspectratio]{confusionmatrix1a.png}
\caption{Confusion matrix of the logistic regression on the test set.}
\label{fig:conf_log}
\end{figure}

\FloatBarrier
\subsection{Question 2}

A perceptron was trained on the same training set and evaluated on the same test set, using the same \texttt{StandardScaler} pipeline. It is the same linear model as the logistic regression, but with no sigmoid: it outputs a hard class decision rather than a probability, and it updates its weights only on misclassified samples.

\FloatBarrier
\subsubsection*{(a) Which model generalized better?}

\begin{table}[H]
\centering
\caption{Test-set performance of the two linear models (spiny = positive class).}
\label{tab:q2a}
\small
\begin{tabular}{lcccc}
\toprule
Model & Accuracy & Precision & Recall & F1 \\
\midrule
Logistic regression & \textbf{0.9860} & \textbf{0.9928} & 0.9786 & \textbf{0.9856} \\
Perceptron          & 0.9614 & 0.9448 & 0.9786 & 0.9614 \\
\bottomrule
\end{tabular}
\end{table}

\textbf{Logistic regression generalizes better.} It wins on accuracy, precision and F1, and ties the perceptron on recall (Table~\ref{tab:q2a}). The recall tie is exact: both models correctly identify 137 of the 140 spiny neurons.

So the whole difference between the two models is in the \emph{false positives}. The perceptron's test confusion matrix is 137 true negatives, 8 false positives, 3 false negatives and 137 true positives. It finds exactly as many spiny neurons as the logistic regression, but mislabels 8 aspiny neurons as spiny against the logistic regression's 1.

The reason is how each model sets its boundary. Logistic regression minimizes a loss over \emph{every} training sample, including correctly classified ones, so the boundary settles where the two classes are best separated on average. The perceptron stops updating on a sample once it is classified correctly, so its final boundary depends on the last few mistakes it corrected, and it has no reason to move away from the aspiny cloud once it clears it.

\FloatBarrier
\subsubsection*{(b) Does the model overfit?}

\begin{table}[H]
\centering
\caption{Train vs.\ test metrics. A positive gap means the model does worse on unseen data.}
\label{tab:q2b}
\small
\begin{tabular}{llccc}
\toprule
Model & Metric & Train & Test & Gap (train $-$ test) \\
\midrule
Logistic regression & Accuracy  & 0.9824 & 0.9860 & $-0.0035$ \\
                    & Precision & 0.9839 & 0.9928 & $-0.0089$ \\
                    & Recall    & 0.9804 & 0.9786 & $+0.0018$ \\
                    & F1        & 0.9821 & 0.9856 & $-0.0035$ \\
\midrule
Perceptron          & Accuracy  & 0.9675 & 0.9614 & $+0.0061$ \\
                    & Precision & 0.9763 & 0.9448 & $+0.0315$ \\
                    & Recall    & 0.9571 & 0.9786 & $-0.0214$ \\
                    & F1        & 0.9666 & 0.9614 & $+0.0052$ \\
\bottomrule
\end{tabular}
\end{table}

\textbf{No, neither model overfits} (Table~\ref{tab:q2b}). Overfitting shows up as training performance clearly above test performance, and that does not happen here.

Logistic regression scores slightly \emph{higher} on the test set than on the training set on three of the four metrics. A model that memorized its training data could not do that. The perceptron's accuracy gap is $+0.0061$, which is under two samples out of 285 and within noise.

The perceptron's precision and recall gaps are worth a closer look, because they point in \emph{opposite} directions: precision drops by $0.0315$ from train to test while recall rises by $0.0214$. Overfitting would make both worse. Instead, the error type flips between the two sets. On the training set it makes more false negatives than false positives (24 vs.\ 13); on the test set it makes more false positives (8 vs.\ 3). The boundary sits in a slightly different place relative to each set, trading precision for recall, which fits the instability described in (a). It does not indicate that the model memorized the training data.

\FloatBarrier
\subsection{Question 3}

A multilayer perceptron was trained to classify the two neuron types. The input is a table of numeric features, so all layers are fully connected. The features were standardized before training with the same \texttt{StandardScaler} fitted on the training set only, and the neuron-type column was not used as an input. All networks use a 2-unit output layer trained with cross-entropy loss.

\FloatBarrier
\subsubsection*{(a) Comparison of two MLP architectures}

The two architectures differ in depth, activation function and regularization (Table~\ref{tab:q3a}):
\begin{itemize}
    \item \textbf{mlp}: one hidden layer, ReLU, no regularization. This is the smaller baseline.
    \item \textbf{mlp-dropout}: two hidden layers, GELU, and dropout of 0.2 after each hidden layer.
\end{itemize}

\begin{table}[H]
\centering
\caption{The two MLP architectures, their training settings, and their accuracy.}
\label{tab:q3a}
\small
\begin{tabular}{lcc}
\toprule
 & \textbf{mlp} & \textbf{mlp-dropout} \\
\midrule
Layer widths            & $41 \rightarrow 32 \rightarrow 2$ & $41 \rightarrow 32 \rightarrow 32 \rightarrow 2$ \\
Activation functions    & ReLU & GELU \\
Regularization          & none & dropout 0.2 \\
Trainable parameters    & 1{,}410 & 2{,}466 \\
Optimizer               & Adam & Adam \\
Learning rate           & 0.001 & 0.001 \\
Epochs                  & 100 & 100 \\
Batch size              & 32 & 32 \\
\midrule
Train accuracy          & 1.0000 & 1.0000 \\
Test accuracy           & 0.9860 & \textbf{0.9930} \\
\bottomrule
\end{tabular}
\end{table}

The parameter counts include weights and biases:
\begin{itemize}
    \item \textbf{mlp}: $41 \cdot 32 + 32 = 1{,}344$ in the hidden layer and $32 \cdot 2 + 2 = 66$ in the output layer, giving 1{,}410.
    \item \textbf{mlp-dropout}: the same 1{,}344 and 66, plus $32 \cdot 32 + 32 = 1{,}056$ for the second hidden layer, giving 2{,}466. Dropout layers have no parameters.
\end{itemize}


The deeper, regularized architecture scores higher on the test set: 0.9930 against 0.9860, which is two fewer errors, using $1.75\times$ the parameters. Both networks reach 1.0000 accuracy on the training set, so both memorize it completely, yet \textbf{mlp-dropout} generalizes better.

This comparison cannot tell us \emph{which} change caused the improvement. The two architectures differ in depth, activation and dropout at the same time, so the gain belongs to the combination. Part (c) varies one hyper-parameter at a time to separate these effects, and Question~4(d) tests dropout on its own.

\FloatBarrier
\subsubsection*{(b) MLP against the logistic regression and perceptron baselines}

\begin{table}[H]
\centering
\caption{Test-set performance of all four models (spiny = positive class).}
\label{tab:q3b}
\small
\begin{tabular}{lcccc}
\toprule
Model & Accuracy & Precision & Recall & F1 \\
\midrule
Logistic regression    & 0.9860 & 0.9928 & 0.9786 & 0.9856 \\
Perceptron             & 0.9614 & 0.9448 & 0.9786 & 0.9614 \\
mlp           & 0.9860 & 0.9857 & 0.9857 & 0.9857 \\
mlp-dropout  & \textbf{0.9930} & \textbf{0.9929} & \textbf{0.9929} & \textbf{0.9929} \\
\bottomrule
\end{tabular}
\end{table}

\textbf{mlp-dropout} is the best model on all four metrics (Table~\ref{tab:q3b}). The margin is small, though. It beats the logistic regression by 0.0070 in accuracy, which is \textbf{two samples} out of 285. The plain \texttt{mlp} has exactly the same accuracy as the logistic regression, with slightly lower precision and slightly higher recall. The perceptron is the only model that clearly falls behind.

The main conclusion is that this task is \textbf{close to linearly separable}. A linear classifier already reaches 98.6\%, and a network with 2{,}466 parameters adds only two more correct samples. For this dataset, the extra capacity of an MLP brings a marginal gain at best.

The models also behave differently on the training set. Both MLPs reach 1.0000 training accuracy, giving train--test gaps of $+0.0140$ (\textbf{mlp}) and $+0.0070$ (\textbf{mlp-dropout}). The logistic regression has a \emph{negative} gap ($-0.0035$). So the networks memorize their training set and the linear models do not, even though their test accuracy is nearly the same. This is analyzed further in Question~4(c).

\FloatBarrier
\subsubsection*{(c) Hyper-parameter sensitivity}

Five hyper-parameters were tested \textbf{one at a time}. In each sweep, one hyper-parameter changes while all the others stay at a fixed base configuration: two hidden layers of 32 units, ReLU, no dropout, Adam with learning rate 0.001, 100 epochs and batch size 32. The values tested were:
\begin{itemize}
    \item learning rate: 0.0001, 0.001, 0.01, 0.1
    \item hidden units per layer: 8, 16, 32, 64, 128
    \item number of hidden layers: 1, 2, 3, 4
    \item dropout: 0, 0.2, 0.5
    \item batch size: 8, 16, 32, 64, 128
\end{itemize}
Sweeping every combination at once would mix the effects together, and no single hyper-parameter's influence could be read from the result.

Every configuration was scored on a \textbf{validation set} made of 10\% of the training set (114 samples), \emph{not} on the test set. If hyper-parameters were chosen by test accuracy, the reported test score would be the best of many tries and would no longer be a fair estimate of generalization.

\begin{figure}[H]
\centering
\includegraphics[width=0.98\linewidth,height=0.36\textheight,keepaspectratio]{sensitivity3c.png}
\caption{Validation accuracy as each hyper-parameter is varied, with all others held fixed. Learning rate, hidden units and batch size are plotted on a log scale.}
\label{fig:sensitivity}
\end{figure}

\begin{table}[H]
\centering
\caption{The learning-rate sweep, the only one with a large effect.}
\label{tab:q3c_lr}
\small
\begin{tabular}{ccl}
\toprule
Learning rate & Validation accuracy & \\
\midrule
0.0001 & 0.9825 & converges too slowly for 100 epochs \\
0.001  & 0.9912 & \\
0.01   & 1.0000 & \\
0.1    & \textbf{0.5088} & training diverges \\
\bottomrule
\end{tabular}
\end{table}

\textbf{Only the learning rate has a large effect} (Figure~\ref{fig:sensitivity}). The other four sweeps barely move:
\begin{itemize}
    \item \textbf{Hidden units per layer} and \textbf{dropout} are completely flat, at 0.9912 for every value.
    \item \textbf{Hidden layers} and \textbf{batch size} each change at only one setting, by one sample. Three hidden layers reach 1.0000 (one more correct), and batch size 64 drops to 0.9825 (one more wrong).
\end{itemize}
With 114 validation samples, one sample is worth 0.0088 accuracy. None of these four sweeps shows a difference bigger than a single neuron.

The misclassified samples explain why the curves are so flat. \textbf{The same validation sample is misclassified in 18 of the 21 runs}, and in most runs it is the only error, so 0.9912 is simply $113/114$. This is \emph{saturation}, not noise. Almost every configuration fits the training set perfectly, and the validation set contains one hard example that almost no configuration gets right. That leaves no room for a hyper-parameter to change the score. The sweeps did produce genuinely different models: the parameter count ranges from 426 (8 hidden units) to 22{,}146 (128 hidden units), and the final training loss ranges from $10^{-5}$ to $8 \times 10^{-3}$. The measurement is what saturates.

The learning rate is different (Table~\ref{tab:q3c_lr}). At 0.0001 the network learns too slowly to finish converging in 100 epochs. At \textbf{0.1} training breaks down completely. The training loss stays at 0.680, close to $\ln 2 \approx 0.693$, which is the loss of a classifier that answers 50/50 for every sample. Training accuracy falls to 0.5083, and the network puts \emph{every} sample in the same class. That is why its accuracy of 0.5088 equals the share of aspiny neurons in the validation set. The steps are too large for gradient descent to settle into a minimum.

\textbf{Conclusion.} This does not mean architecture never matters. The task is nearly linearly separable, so every reasonable architecture reaches the ceiling, and the only hyper-parameter that can do real damage is the one that breaks the optimization itself.

\FloatBarrier
\subsection{Question 4}

A network with one hidden layer (32 units, ReLU) was trained with full-batch \textbf{Gradient Descent (GD)}. Each epoch computes the gradient over the whole training set and makes exactly one parameter update. 10\% of the training set (114 neurons) was held out as a validation set, leaving 1{,}025 neurons for training. Training used the plain SGD optimizer (no momentum, not Adam) with learning rate 0.01 for 300 epochs.

For part (b) the same network was trained again with \textbf{Stochastic Gradient Descent (SGD)} with batch size 10. Both runs use the same data split, the same initial weights (same random seed), the same optimizer and the same learning rate. The only difference between them is the batch size. Figure~\ref{fig:loss_gd_sgd} shows the training and validation loss of both runs.

\begin{figure}[H]
\centering
\includegraphics[width=0.98\linewidth,height=0.36\textheight,keepaspectratio]{loss_gd_sgd4.png}
\caption{Training and validation cross-entropy loss per epoch. Left: full-batch GD. Right: SGD with batch size 10.}
\label{fig:loss_gd_sgd}
\end{figure}

\FloatBarrier
\subsubsection*{(a) Error analysis of the GD-trained model}

\begin{figure}[H]
\centering
\includegraphics[width=0.45\linewidth,height=0.36\textheight,keepaspectratio]{confusionmatrix4a.png}
\caption{Confusion matrix of the GD-trained MLP on the test set.}
\label{fig:conf_gd}
\end{figure}

The GD-trained model reaches accuracy $=0.9333$, precision $=0.8954$, recall $=0.9786$ and F1 $=0.9352$. Its confusion matrix (Figure~\ref{fig:conf_gd}) gives: true negatives 129, false positives 16, false negatives 3, true positives 137.

The 19 errors are strongly one-sided:
\begin{itemize}
    \item 16 false positives --- aspiny neurons classified as spiny
    \item 3 false negatives --- spiny neurons classified as aspiny
\end{itemize}
So precision (0.8954) is far below recall (0.9786), the opposite of the logistic regression in Question~1. This model is too quick to assign the spiny label.

This is the weakest model in Part~C, worse even than the perceptron. The reason is not the architecture but the training method. Full-batch GD makes \textbf{one parameter update per epoch}, so 300 epochs means only 300 updates. SGD with batch size 10 makes 103 updates per epoch, or 30{,}900 over the same 300 epochs. Figure~\ref{fig:loss_gd_sgd} shows that the GD loss is still falling at epoch 300: the model simply has not finished training.

\FloatBarrier
\subsubsection*{(b) Comparing the GD and SGD loss curves}

\begin{table}[H]
\centering
\caption{Loss at the end of training (epoch 300).}
\label{tab:q4b}
\small
\begin{tabular}{lccc}
\toprule
Method & Parameter updates & Final train loss & Final validation loss \\
\midrule
GD (full batch)  & 300    & 0.2370 & 0.2668 \\
SGD (batch = 10) & 30{,}900 & 0.0037 & 0.0583 \\
\bottomrule
\end{tabular}
\end{table}

\textbf{SGD converges much faster per epoch} (Table~\ref{tab:q4b}). With the same learning rate and the same 300 epochs, it reaches a training loss about 60 times lower than GD. Each SGD step is actually \emph{less} accurate than a GD step, since it estimates the gradient from 10 samples instead of all 1{,}025. But SGD takes 103 steps for every one GD step. The SGD training loss drops below 0.1 by epoch 15. The GD loss falls smoothly from 0.71 and is still going down at the end (0.2593 at epoch 250, 0.2370 at epoch 300).

Both curves in Figure~\ref{fig:loss_gd_sgd} look smooth, which may seem surprising for SGD. Individual mini-batch steps are noisy, but each point on the curve is the \emph{average} loss over all 103 mini-batches in that epoch, and the averaging hides the noise.

The curves differ most in the second half of training. For GD, the validation loss stays close to the training loss the whole time. For SGD, the validation loss stops improving around epoch 112 and then slowly rises, while the training loss keeps falling towards zero. This is the classic sign of overfitting, analyzed in part (c).

\FloatBarrier
\subsubsection*{(c) Does the MLP overfit?}

\begin{table}[H]
\centering
\caption{Train vs.\ test accuracy of the two MLPs. Gap = train $-$ test.}
\label{tab:q4c}
\small
\begin{tabular}{lcccl}
\toprule
Model & Train accuracy & Test accuracy & Gap & Diagnosis \\
\midrule
GD  & 0.9200 & 0.9333 & $-0.0133$ & underfitting \\
SGD & 1.0000 & 0.9860 & $+0.0140$ & overfitting \\
\bottomrule
\end{tabular}
\end{table}

The answer depends on how the network was trained (Table~\ref{tab:q4c}).

\textbf{The GD-trained model does not overfit; it underfits.} It scores \emph{higher} on the test set than on the training set. Its training and validation losses are both high (0.2370 and 0.2668), stay close together, and are still falling at the last epoch. A model that has not yet learned its training set cannot be memorizing it.

\textbf{The SGD-trained model does overfit}, and the loss curve shows when it starts. The validation loss reaches its lowest value, \textbf{0.0501, at epoch 112}. After that it rises to 0.0583 by epoch 300, while the training loss keeps falling to 0.0037. By the end, the validation loss is about 16 times the training loss. So every epoch after 112 improves the fit on the training data while making the model slightly worse on unseen data. The test metrics agree: 1.0000 accuracy on the training set against 0.9860 on the test set.

The overfitting is mild. An accuracy gap of $+0.0140$ is 4 neurons out of 285, and the model still generalizes well. Stopping training at around epoch 112, when the validation loss is lowest (early stopping), would have avoided most of it.

\FloatBarrier
\subsubsection*{(d) Adding 50\% dropout before the classification layer}

A dropout layer with rate 0.5 was added between the hidden layer and the output layer, so it sits directly before the classification layer. Everything else is the same as the SGD run: batch size 10, learning rate 0.01, 300 epochs. The comparison is therefore against the SGD model, which is the model that overfitted in part (c).

\begin{figure}[H]
\centering
\includegraphics[width=0.70\linewidth,height=0.36\textheight,keepaspectratio]{loss_dropout4d.png}
\caption{Training and validation loss of the SGD model with 50\% dropout before the output layer.}
\label{fig:loss_dropout}
\end{figure}

\begin{table}[H]
\centering
\caption{Effect of dropout at the end of training.}
\label{tab:q4d}
\small
\begin{tabular}{lcccc}
\toprule
Model & Train acc. & Test acc. & Gap & Lowest validation loss \\
\midrule
SGD, no dropout  & 1.0000 & 0.9860 & $+0.0140$ & 0.0501 (epoch 112) \\
SGD, dropout 0.5 & 0.9990 & \textbf{0.9930} & $\mathbf{+0.0060}$ & \textbf{0.0381} (epoch 121) \\
\bottomrule
\end{tabular}
\end{table}

\textbf{Yes, dropout improved generalization} (Table~\ref{tab:q4d}), on three measures:
\begin{enumerate}
    \item \textbf{Test accuracy increased} from 0.9860 to 0.9930, two more neurons classified correctly.
    \item \textbf{The train--test gap was more than halved}, from $+0.0140$ to $+0.0060$. Training accuracy is no longer a perfect 1.0000, so the network stopped memorizing the training set completely.
    \item \textbf{The lowest validation loss dropped by 24\%}, from 0.0501 to 0.0381.
\end{enumerate}

Figure~\ref{fig:loss_dropout} needs careful reading. The training loss is recorded \emph{with dropout switched on}, while the validation loss is computed with dropout switched off. This explains two things in the plot. The training curve is noisy, because each step drops a different random half of the hidden units. And the training loss stays \emph{above} the validation loss for the first $\sim$120 epochs, because the network is handicapped during training but not during validation. The loss values are therefore not a like-for-like comparison. The accuracy gap in Table~\ref{tab:q4d}, where both sets are evaluated with dropout off, is the fair measure.

Dropout reduced overfitting but did not remove it. After epoch 121 the validation loss rises again (to 0.0537 by epoch 300), just more slowly than without dropout.

Dropout works because, at every training step, it switches off a random half of the hidden units. The output layer can never rely on one particular unit or combination of units, so the network has to spread what it learns across many of them. The resulting model is less sensitive to the small differences between the training neurons and new ones.

\FloatBarrier
\subsubsection*{(e) Error analysis}

The SGD-trained network misclassifies 4 of the 285 test neurons. We checked each of them in two ways.

\paragraph{Do they resemble the other class?}
The features were ranked by how well they separate spiny from aspiny neurons in the training set, using Cohen's $d$ (the difference between the class means divided by the pooled standard deviation). The 10 most separating features all describe the \emph{shape of the action potential}, for example the upstroke/downstroke ratio, the downstroke velocity and the peak voltage. For each test neuron we counted on how many of these 10 features it lies closer to the mean of the \emph{other} class than to its own (Table~\ref{tab:q4e}).

\begin{table}[H]
\centering
\caption{Misclassified test neurons, compared on the 10 features that best separate the two classes.}
\label{tab:q4e}
\small
\begin{tabular}{cccc}
\toprule
Test index & True class & $P(\text{spiny})$ & Closer to the other class on \\
\midrule
176 & spiny  & 0.174 & 5 of 10 features \\
187 & aspiny & 0.632 & 4 of 10 features \\
245 & spiny  & 0.432 & 4 of 10 features \\
257 & aspiny & 0.925 & \textbf{7 of 10 features} \\
\midrule
\multicolumn{3}{l}{281 correctly classified neurons} & median 1 of 10 \\
\bottomrule
\end{tabular}
\end{table}

A correctly classified neuron is closer to the other class on a median of only \textbf{1} of these features. Every misclassified neuron is closer to the other class on \textbf{4 to 7}. On exactly the features the model relies on most, these neurons look like the class it assigned them to.

\paragraph{Are they unusual for their own class?}
Each misclassified neuron also lies far from its own class mean on at least one feature. The most extreme is neuron 257, whose \texttt{sag} is \textbf{3.72 standard deviations} above the aspiny mean. However, the features on which these neurons are most unusual (\texttt{sag}, \texttt{input\_resistance\_mohm}, \texttt{threshold\_v}) separate the classes only weakly: \texttt{sag} ranks 30th of the 41 features and \texttt{threshold\_v} ranks 39th. Being unusual on these features does not by itself explain the errors. The resemblance on the strongly separating features does.

\paragraph{How confident was the model?}
Neurons 187 ($P=0.632$) and 245 ($P=0.432$) are close to the 0.5 decision boundary, so they are borderline cases. Neurons 176 ($P=0.174$) and 257 ($P=0.925$) are \textbf{confidently wrong}. For these two, the electrical features point clearly to the other class. Neuron 257 matches the spiny profile on 7 of the 10 most informative features.

\paragraph{Conclusion.}
The errors are not random. They reflect a real overlap between the two cell types in this feature space. A few neurons have the electrical signature of the other type, and no classifier using these features alone could be expected to classify them correctly. This could be natural biological variation, or these neurons could have incorrect labels in the dataset.

\FloatBarrier
\subsection{Conclusion}

\begin{table}[H]
\centering
\caption{Test performance of all models in Part C.}
\label{tab:conclusion}
\small
\begin{tabular}{lcc}
\toprule
Model & Test accuracy & Test F1 \\
\midrule
Logistic regression                      & 0.9860 & 0.9856 \\
Perceptron                               & 0.9614 & 0.9614 \\
MLP, 1 hidden layer (Adam)               & 0.9860 & 0.9857 \\
MLP, 2 hidden layers + dropout 0.2 (Adam) & \textbf{0.9930} & \textbf{0.9929} \\
MLP, GD                                  & 0.9333 & 0.9352 \\
MLP, SGD (batch 10)                      & 0.9860 & 0.9857 \\
MLP, SGD + dropout 0.5                   & \textbf{0.9930} & \textbf{0.9929} \\
\bottomrule
\end{tabular}
\end{table}

\begin{itemize}
    \item \textbf{The task is almost linearly separable.} Logistic regression already reaches 98.6\%, and the best MLP is only two neurons better.
    \item \textbf{Training mattered more than architecture.} The same network got 0.9333 with GD and 0.9860 with SGD, because GD made far fewer updates. In the sweep, only the learning rate had a real effect.
    \item \textbf{The linear models did not overfit; the MLPs did, mildly.} The MLPs reached 100\% training accuracy, and adding dropout halved the train--test gap and gave the best test accuracy.
    \item \textbf{The few remaining errors are neurons that look like the other class} on the most informative features, so these features alone probably can't classify them correctly.
\end{itemize}


\end{document}
